% Part: first-order-logic
% Chapter: completeness
% Section: henkin-expansion

\documentclass[../../../include/open-logic-section]{subfiles}

\begin{document}

\olfileid[tr]{fol}{com}{hen}

\olsection{Henkin Genişlemesi}

\begin{explain}
Tamlık teoremini ispatlamaktaki güçlüğün bir bölümü, tam ve tutarlı
bir~$\Gamma$ kümesinden kurduğumuz modelin~$\Gamma$ içindeki bütün
niceleyicili !!{formula}s ifadelerini doğru kılması gereğidir. Bunu
güvence altına almak için Leon Henkin'e dayanan bir hile kullanırız.
Hile özünde, dili sonsuz sayıda !!{constant}s ile genişletmekten ve her
$!A(x)$ !!{formula} ifadesi için---burada bir !!{variable} serbesttir---
\iftag{prvEx}
      {$\lexists[x][!A(x)]\lif !A(c)$}
      {$\lnot\lforall[x][!A(x)]\lif\lnot !A(c)$}
biçiminde bir formül eklemekten oluşur; burada $c$ yeni
!!{constant}s arasındadır. $\Gamma$'yı sağlayan !!{structure}
kurulduğunda bu, her
\iftag{prvEx}
{doğru varlıksal tümcenin bir tanığının}
{yanlış tümel tümcenin bir karşı örneğinin}
yeni sabitler arasında bulunmasını güvence altına alır.
\end{explain}

\begin{prop}
\ollabel{prop:lang-exp}
$\Gamma$,~$\Lang L$ içinde tutarlıysa ve $\Lang L'$,
$\Lang L$ diline !!a{denumerable} yeni !!{constant}s kümesi
$\Obj d_0$, $\Obj d_1$, \dots\ eklenerek elde edilmişse,
$\Gamma$~$\Lang L'$ içinde de tutarlıdır.
\end{prop}

\begin{defn}[Doygun küme]
Bir $\Lang L$ dilindeki !!{formula}s kümesi~$\Gamma$, her
$!A(x)\in\Frm[L]$ !!{formula} ifadesi için---burada bir !!{variable}~$x$
serbesttir---
öyle bir !!a{constant}~$c\in\Lang L$ varsa ki
\iftag{prvEx}
      {$\lexists[x][!A(x)]\lif !A(c)\in\Gamma$}
      {$\lnot\lforall[x][!A(x)]\lif\lnot !A(c)\in\Gamma$}
olsun, \emph{doygun}dur.
\end{defn}

Aşağıdaki tanım bir sonraki teoremin ispatında kullanılacaktır.

\begin{defn}
\ollabel{defn:henkin-exp}
$\Lang L'$,~\olref{prop:lang-exp} önermesindeki gibi olsun.
$\Lang L'$ dilinde bir değişkenin ($x_i$) serbest geçtiği bütün
$!A_i(x_i)$ !!{formula}s ifadelerinin $!A_0(x_0)$, $!A_1(x_1)$, \dots\
numaralandırmasını sabitleyelim. !!{sentence}s~$!D_n$ ifadelerini~$n$
üzerinde tümevarımla tanımlarız.

$c_0$,~$\Lang L$ diline eklediğimiz $\Obj d_i$ simgeleri arasından
$!A_0(x_0)$ içinde geçmeyen ilk !!{constant} olsun. $!D_0$, \dots,
$!D_{n-1}$ ifadelerinin tanımlanmış olduğunu varsayarak $c_n$'yi, yeni
$\Obj d_i$ !!{constant}s arasından ne $!D_0$, \dots,~$!D_{n-1}$
ifadelerinde ne de $!A_n(x_n)$ içinde geçen ilk sabit olarak seçelim.

Şimdi $!D_n$ şu !!{formula} olsun: \iftag{prvEx}
{$\lexists[x_n][!A_n(x_n)]\lif !A_n(c_n)$}
{$\lnot\lforall[x_n][!A_n(x_n)]\lif\lnot !A_n(c_n)$}.
\end{defn}

\begin{lem}
\ollabel{lem:henkin}
Her tutarlı~$\Gamma$ kümesi, doygun ve tutarlı bir~$\Gamma'$ kümesine
genişletilebilir.
\end{lem}

\begin{proof}
Bir~$\Lang L$ dilindeki tutarlı tümceler kümesi~$\Gamma$ verilsin.
!!a{denumerable} yeni !!{constant}s kümesi ekleyerek
dili~$\Lang L'$ diline genişletelim. \olref{prop:lang-exp} uyarınca
$\Gamma$ bu daha zengin dilde de tutarlıdır. Ayrıca $!D_i$ ifadeleri
\olref{defn:henkin-exp} tanımındaki gibi olsun. Şöyle tanımlayalım:
\begin{align*}
\Gamma_0 & = \Gamma \\
\Gamma_{n+1} & = \Gamma_n \cup \{!D_n \}.
\end{align*}
Başka bir deyişle $\Gamma_{n+1}=\Gamma\cup\{!D_0,\dots,!D_n\}$ ve
$\Gamma'=\bigcup_n\Gamma_n$ olsun. $\Gamma'$ açıkça doygundur.

$\Gamma'$ tutarsız olsaydı bir~$n$ için $\Gamma_n$ tutarsız olurdu
(Alıştırma: nedenini açıklayın). Dolayısıyla $\Gamma'$ kümesinin
tutarlı olduğunu göstermek için her~$\Gamma_n$ kümesinin tutarlı
olduğunu~$n$ üzerinde tümevarımla göstermek yeterlidir.

Tümevarımın temel durumu, teoremin varsayımı olan $\Gamma_0=\Gamma$
kümesinin tutarlı olduğu savıdır. Tümevarım adımında $\Gamma_n$
tutarlı olduğu halde $\Gamma_{n+1}=\Gamma_n\cup\{!D_n\}$ kümesinin
tutarsız olduğunu varsayalım. $!D_n$ ifadesinin
\iftag{prvEx}
{$\lexists[x_n][!A_n(x_n)]\lif !A_n(c_n)$}
{$\lnot\lforall[x_n][!A_n(x_n)]\lif\lnot !A_n(c_n)$}
olduğunu hatırlayalım; burada $!A_n(x_n)$,~$\Lang L'$ dilinde yalnızca
$x_n$ değişkeninin serbest olduğu !!a{formula} ifadesidir.
$c_n$'yi seçme biçimimizden ötürü (bkz. \olref{defn:henkin-exp})
$c_n$, ne $!A_n(x_n)$ içinde ne de~$\Gamma_n$ içinde geçer.

$\Gamma_n\cup\{!D_n\}$ tutarsızsa $\Gamma_n\Proves\lnot !D_n$ ve
dolayısıyla aşağıdakilerin ikisi de geçerlidir:
\iftag{prvEx}{
\[
\Gamma_n \Proves \lexists[x_n][!A_n(x_n)]
\qquad
\Gamma_n \Proves \lnot !A_n(c_n)
\]}{
\[
\Gamma_n \Proves \lnot\lforall[x_n][!A_n(x_n)]
\qquad
\Gamma_n \Proves !A_n(c_n).
\]}
$c_n$,~$\Gamma_n$ ya da~$!A_n(x_n)$ içinde geçmediğinden
\tagrefs{prfAX/{fol:axd:qpr:thm:strong-generalization},prfSC/{fol:seq:qpr:thm:strong-generalization},prfND/{fol:ntd:qpr:thm:strong-generalization},prfTab/{fol:tab:qpr:thm:strong-generalization}}
uygulanabilir.
\iftag{prvEx}
{$\Gamma_n\Proves\lnot !A_n(c_n)$}
{$\Gamma_n\Proves !A_n(c_n)$}
ifadesinden
\iftag{prvEx}
{$\Gamma_n\Proves\lforall[x_n][\lnot !A_n(x_n)]$}
{$\Gamma_n\Proves\lforall[x_n][!A_n(x_n)]$}
elde ederiz. Böylece aynı anda
\iftag{prvEx}
{$\Gamma_n\Proves\lexists[x_n][!A_n(x_n)]$ ve
$\Gamma_n\Proves\lforall[x_n][\lnot !A_n(x_n)]$}
{$\Gamma_n\Proves\lnot\lforall[x_n][!A_n(x_n)]$ ve
$\Gamma_n\Proves\lforall[x_n][!A_n(x_n)]$}
olur; dolayısıyla $\Gamma_n$ kümesinin kendisi tutarsızdır.
\iftag{prvEx}
{(Şuna dikkat edin:
\iftag{prvAll}
{$\lforall[x_n][\lnot !A_n(x_n)]\Proves
  \lnot\lexists[x_n][!A_n(x_n)]$}
{$\lforall[x_n][\lnot !A_n]$ ifadesi
  $\lnot\lexists[x_n][\lnot\lnot !A_n(x_n)]$ olarak tanımlanır ve
  $\lnot\lexists[x_n][\lnot\lnot !A_n(x_n)]\Proves
  \lnot\lexists[x_n][!A_n(x_n)]$}.)}{}
Bu bir çelişkidir: $\Gamma_n$'nin tutarlı olduğu varsayılmıştı.
Öyleyse $\Gamma_n\cup\{!D_n\}$ tutarlıdır.
\end{proof}

\begin{explain}
Şimdi doygun olan \emph{tam} ve tutarlı kümelerin şu özelliğe sahip
olduğunu göstereceğiz: \iftag{prvAll}{tümel niceleyicili bir
  !!{sentence} ifadesini ancak ve ancak onun bütün örneklerini
  içerdiğinde içerir}{}\iftag{defAll,defEx}{}{ ve }
\iftag{prvEx}{varlıksal niceleyicili bir !!{sentence} ifadesini ancak
  ve ancak onun en az bir örneğini içerdiğinde içerir}{}. Bunu, tam,
tutarlı ve doygun bir kümeden üreteceğimiz !!{structure} ifadesinin
kümedeki niceleyicili bütün tümceleri doğru kıldığını göstermek için
kullanacağız.
\end{explain}

\begin{prop}\ollabel{prop:saturated-instances}
$\Gamma$'nın tam, tutarlı ve doygun olduğunu varsayalım.
\begin{tagenumerate}{prvEx,prvAll}
\tagitem{prvEx}{$\lexists[x][!A(x)]\in\Gamma$ ancak ve ancak en az bir
  kapalı~$t$ terimi için $!A(t)\in\Gamma$ ise.}{}
\tagitem{prvAll}{$\lforall[x][!A(x)]\in\Gamma$ ancak ve ancak bütün
  kapalı~$t$ terimleri için $!A(t)\in\Gamma$ ise.}{}
\end{tagenumerate}
\end{prop}

\begin{proof}
  \begin{tagenumerate}{prvEx,prvAll}
  \tagitem{prvEx}{%
    \iftag{probEx}{Alıştırma.}{Önce $\lexists[x][!A(x)]\in\Gamma$
      olduğunu varsayalım. $\Gamma$ doygun olduğundan bir
      !!{constant}~$c$ için $(\lexists[x][!A(x)]\lif !A(c))\in\Gamma$.
      \tagrefs{prfAX/{fol:axd:ppr:prop:provability-lif},prfSC/{fol:seq:ppr:prop:provability-lif},prfND/{fol:ntd:ppr:prop:provability-lif},prfTab/{fol:tab:ppr:prop:provability-lif}}
      sonucunun (1). maddesi ile
      \olref[ccs]{prop:ccs}\olref[ccs]{prop:ccs-prov-in} uyarınca
      $!A(c)\in\Gamma$.

    Öteki yönde doygunluk gerekmez: $!A(t)\in\Gamma$ olduğunu
    varsayalım. \tagrefs{prfAX/{fol:axd:qpr:prop:provability-quantifiers},prfSC/{fol:seq:qpr:prop:provability-quantifiers},prfND/{fol:ntd:qpr:prop:provability-quantifiers},prfTab/{fol:tab:qpr:prop:provability-quantifiers}}
    sonucunun (1). maddesine göre $\Gamma\Proves\lexists[x][!A(x)]$.
    \olref[ccs]{prop:ccs}\olref[ccs]{prop:ccs-prov-in} uyarınca
    $\lexists[x][!A(x)]\in\Gamma$.}}{}
  \tagitem{prvAll}{%
    \iftag{probAll}{Alıştırma.}{Bütün kapalı~$t$ terimleri için
      $!A(t)\in\Gamma$ olduğunu varsayalım. Olmayana ergiyle
      $\lforall[x][!A(x)]\notin\Gamma$ olsun. $\Gamma$ tam olduğundan
      $\lnot\lforall[x][!A(x)]\in\Gamma$. Doygunluk uyarınca bir
      !!{constant}~$c$ için \iftag{prvEx}{$(\lexists[x][\lnot !A(x)]
        \lif\lnot !A(c))\in\Gamma$}{$(\lnot\lforall[x][!A(x)]
        \lif\lnot !A(c))\in\Gamma$}. Varsayıma göre $c$ kapalı bir
      terim olduğundan $!A(c)\in\Gamma$. Fakat bu,~$\Gamma$'yı
      tutarsız kılar. (Alıştırma: Aşağıdakinin tutarsız olduğunu
      gösteren !!{derivation} ifadesini verin:
      \[
      \lnot \lforall[x][!A(x)],
      \iftag{prvEx}{\lexists[x][\lnot !A(x)] \lif \lnot
        !A(c)}{\lnot\lforall[x][!A(x)] \lif \lnot
        !A(c)}, !A(c).
      \])

      Ters yönde doygunluk gerekmez: $\lforall[x][!A(x)]\in\Gamma$
      olduğunu varsayalım.
      \tagrefs{prfAX/{fol:axd:qpr:prop:provability-quantifiers},prfSC/{fol:seq:qpr:prop:provability-quantifiers},prfND/{fol:ntd:qpr:prop:provability-quantifiers},prfTab/{fol:tab:qpr:prop:provability-quantifiers}}
      sonucunun (2). maddesine göre $\Gamma\Proves !A(t)$.
      \olref[ccs]{prop:ccs} uyarınca $!A(t)\in\Gamma$ elde ederiz.}}{}
  \end{tagenumerate}
\end{proof}

\end{document}
